\chapter{Agriculturals and Softs}\label{ch:m3:agriculturals-and-softs}

At noon Eastern time on 30 June 2026 the US Department of Agriculture published its estimate of
the acreage American farmers had planted that spring: corn down 3\% on the year, soybeans up 5\%.
For the traders watching, the minutes after noon were the most important of the season: the
number told them how much corn would exist in the autumn, and the futures repriced it before most
farmers had read the release. Agricultural markets run on calendars that nature sets and
governments measure. This chapter covers the \omterm{def:m3:agriculturals-and-softs:cropyear}{crop year} and its spreads, the reports that move the
market, the weather, the daily limits that can lock a trader in, and the weekly report that shows
who holds what.

\section{Crop calendars and the old-crop, new-crop spread}

\begin{definition}[Crop year, old crop, new crop]\label{def:m3:agriculturals-and-softs:cropyear}
The \emph{crop year}\index{crop year} (marketing year) of a crop is the twelve-month period over
which one harvest is sold and used; for US corn it runs from 1 September to 31 August.
\emph{Old crop}\index{old crop} is grain from the harvest already in store; \emph{new
crop}\index{new crop} is grain from the harvest to come. Futures delivering before the new harvest
price old crop, those after it new crop.
\end{definition}

\begin{omfigure}
\centering
\begin{tikzpicture}[font=\small,x=0.9cm]
\draw[thick] (0,0) -- (12,0);
\foreach \i/\m in {0/Jan,1/Feb,2/Mar,3/Apr,4/May,5/Jun,6/Jul,7/Aug,8/Sep,9/Oct,10/Nov,11/Dec}{
  \draw (\i,0.1) -- (\i,-0.1); \node[anchor=north,font=\footnotesize] at (\i+0.5,-0.1) {\m};}
\fill[omProp!35] (3.5,0.3) rectangle (5.5,0.8); \node[font=\footnotesize] at (4.5,0.55) {planting};
\fill[omThm!35] (6.4,0.3) rectangle (7.3,0.8); \node[font=\footnotesize,anchor=south west] at (6.4,0.85) {pollination};
\fill[omDef!35] (8.5,0.3) rectangle (10.8,0.8); \node[font=\footnotesize] at (9.65,0.55) {harvest};
\draw[{Stealth[length=2mm]}-{Stealth[length=2mm]},thick] (8,-0.9) -- (12,-0.9) node[midway,below,font=\footnotesize] {new crop year};
\draw[thick,-{Stealth[length=2mm]}] (0,-0.9) -- (7.9,-0.9) node[midway,below,font=\footnotesize] {old crop sold from store};
\node[font=\footnotesize,anchor=south east] at (5.95,1.1) {30 June: acreage};
\draw[gray,dashed] (5.97,0.2) -- (5.97,1.1);
\end{tikzpicture}

\omcaption{The US corn year, schematic: planting in spring, pollination in July (the weather-sensitive
weeks), harvest from autumn; the marketing year begins on 1 September. The acreage report of late June
is the season's first hard number on the size of the \omterm{def:m3:agriculturals-and-softs:cropyear}{new crop}.}\label{fig:m3:agriculturals-and-softs:calendar}
\end{omfigure}

The spread between the last old-crop contract and the first new-crop contract (for corn, July against
December) is a bet on the size of the new harvest against the stocks left from the old. A short \omterm{def:m3:agriculturals-and-softs:cropyear}{old
crop} and a big new one put the July far above December; a glut of old grain and a failing \omterm{def:m3:agriculturals-and-softs:cropyear}{new crop}
invert it.

\section{Government reports}

\begin{definition}[WASDE report, stocks-to-use ratio]\label{def:m3:agriculturals-and-softs:wasde}
The \emph{WASDE report}\index{WASDE report} (World Agricultural Supply and Demand Estimates) is the US
Department of Agriculture's monthly balance sheet of supply, use and ending stocks for the major crops,
in the United States and the world. The \emph{stocks-to-use ratio}\index{stocks-to-use ratio} is the
ending stocks of a marketing year divided by that year's total use: the buffer the market will carry
into the next harvest.
\end{definition}

The price of a storable crop is set by how scarce it will be at the end of the year: a low
\omterm{def:m3:agriculturals-and-softs:wasde}{stocks-to-use ratio} means a thin buffer and a price that must ration use; a high one means storage
must be paid for. Reports that revise the balance sheet move the price immediately.

\begin{dated}{2026-09}{Report times}\label{dat:m3:agriculturals-and-softs:reports}
WASDE is released at 12:00 Eastern time, usually between the 8th and the 12th of each month. The
Acreage report of the National Agricultural Statistics Service is released at noon Eastern at the end of
June; on 30 June 2026 it put corn planted acreage 3\% below 2025 and soybeans 5\% above.
\end{dated}

\begin{example}[A stocks-to-use surprise]\label{ex:m3:agriculturals-and-softs:stu}
A balance sheet with 15.0 billion bushels of use and 1.8 billion of ending stocks has a \omterm{def:m3:agriculturals-and-softs:wasde}{stocks-to-use
ratio} of 12.0\%. A report that cuts the crop by 300 million bushels, with use unchanged, cuts ending
stocks to 1.5 billion: 10.0\%. The price must rise until use falls enough to rebuild the buffer the
market wants to hold.
\end{example}

\section{Weather}

For a growing crop, weather is the dominant risk, and it is concentrated in a few weeks: for corn,
pollination in July, when heat and drought cut yields most. Traders buy that risk ahead of the
weeks and sell it afterwards, a pattern called the weather premium; forecasts from the main
numerical models move prices during the season as reliably as reports do.

\section{Daily limits}

Grain futures have daily price limits: a contract cannot trade more than a set amount above or
below the previous settlement. The limit is a price limit of the kind One Quant Book 1, chapter 12,
described for shares.

\begin{definition}[Limit-locked market, expanded limit]\label{def:m3:agriculturals-and-softs:limit}
A \emph{limit-locked market}\index{limit-locked market} is one whose price has reached the daily limit
with orders on one side only: no trade can happen at a price beyond it, so positions cannot be closed.
An \emph{expanded limit}\index{expanded limit} is the wider limit that applies on the day after a
contract settles at its limit.
\end{definition}

\begin{dated}{2026-09}{Grain limits}\label{dat:m3:agriculturals-and-softs:limits}
CBOT corn futures: 5{,}000 bushels, quoted in cents per bushel. Their limits are reset every six months:
7\% of the average settlement of a reference contract over 45 trading days, rounded to the nearest
5 cents, with a floor of 20 cents; the \omterm{def:m3:agriculturals-and-softs:limit}{expanded limit} is 1.5 times that, rounded up. At an average of
\$4.50 the limit is 30 cents and the \omterm{def:m3:agriculturals-and-softs:limit}{expanded limit} 45 cents. For trade date 8 September 2026 CME
Group listed the corn limit at 30 cents (60 cents for calendar spreads).
\end{dated}

\begin{proposition}[Margin to survive a lock]\label{prop:m3:agriculturals-and-softs:lock}
If news moves the fair price of a contract by $\Delta$ beyond its limit $L$, with \omterm{def:m3:agriculturals-and-softs:limit}{expanded limit} $L_e$,
a holder on the wrong side cannot exit until the price reaches the fair value: it loses $L$ on the
first day and $L_e$ on each further locked day, and needs margin for every locked day before it can
trade. With $L < |\Delta| \le L + L_e$ the lock lasts two days.
\end{proposition}

\begin{proof}
Each locked day the settlement moves by the day's limit and no trade is possible at the fair price;
the lock ends on the first day the remaining gap is within the limit.
\end{proof}

\begin{omfigure}
\begin{tikzpicture}
\begin{axis}[width=11cm,height=4.4cm,
  xlabel={trading day},ylabel={\$ per bushel},
  tick label style={font=\small},label style={font=\small},
  xmin=-0.3,xmax=5.3,ymin=3.5,ymax=4.7,grid=major,grid style={gray!25},xtick={0,1,2,3,4,5},
  legend columns=2,legend style={font=\footnotesize,at={(0.5,-0.33)},anchor=north,draw=none,fill=none,/tikz/every even column/.append style={column sep=8pt}},
  table/col sep=comma]
\addplot[omThm,very thick,mark=*,mark size=1.6pt] table[x=day,y=settle]{figdata/markets-3/09-agriculturals-and-softs/limits.csv};
\addplot[omDef,thick,dashed,const plot] table[x=day,y=fair]{figdata/markets-3/09-agriculturals-and-softs/limits.csv};
\node[font=\footnotesize,anchor=west] at (axis cs:1.05,4.23) {limit down 30 cents, locked};
\node[font=\footnotesize,anchor=west] at (axis cs:2.05,3.83) {expanded limit 45 cents, locked};
\legend{settlement,fair price after the report}
\end{axis}
\end{tikzpicture}

\omcaption{Corn after a report that cuts its fair value from \$4.50 to \$3.70 a bushel: two locked
days at the limit and the \omterm{def:m3:agriculturals-and-softs:limit}{expanded limit}, then trading at the fair price. Illustrative; limits by the
rule of the dated box. Data: the chapter's tutorial.}\label{fig:m3:agriculturals-and-softs:limits}
\end{omfigure}

\section{Positioning reports}

\begin{definition}[Commitments of Traders report, managed money]\label{def:m3:agriculturals-and-softs:cot}
The \emph{Commitments of Traders report}\index{Commitments of Traders report} (COT) is the weekly
publication by the US Commodity Futures Trading Commission of the aggregate long and short positions
of large traders in each futures market, by category, as of Tuesday. In its disaggregated form
\emph{managed money}\index{managed money} is the category of commodity trading advisors, commodity
pool operators and funds that trade on behalf of clients, as distinct from producers and merchants,
swap dealers and other reportable traders.
\end{definition}

\begin{dated}{2026-09}{The COT schedule}\label{dat:m3:agriculturals-and-softs:cot}
The Commission releases the report every Friday at 15:30 Eastern time, with positions as of the
previous Tuesday.
\end{dated}

\begin{omfigure}
\begin{tikzpicture}
\begin{axis}[width=11cm,height=4.8cm,axis y line*=left,
  ylabel={managed money net, \% of OI},
  tick label style={font=\small},label style={font=\small},
  xmin=2016,xmax=2026.8,ymin=-30,ymax=30,grid=major,grid style={gray!25},
  xtick={2016,2018,2020,2022,2024,2026},xticklabel style={/pgf/number format/1000 sep={}},
  table/col sep=comma]
\addplot[omDef,very thick] table[x=t,y=mm]{figdata/markets-3/09-agriculturals-and-softs/mm.csv};\label{plot:m3:mm}
\draw[gray] (axis cs:2016,0) -- (axis cs:2026.8,0);
\end{axis}
\begin{axis}[width=11cm,height=4.8cm,axis y line*=right,axis x line=none,
  ylabel={maize, \$/t},ylabel near ticks,
  tick label style={font=\small},label style={font=\small},
  xmin=2016,xmax=2026.8,ymin=100,ymax=400,
  legend columns=2,legend style={font=\footnotesize,at={(0.5,-0.18)},anchor=north,draw=none,fill=none,/tikz/every even column/.append style={column sep=8pt}},
  table/col sep=comma]
\addlegendimage{/pgfplots/refstyle=plot:m3:mm}\addlegendentry{managed money net (left)}
\addplot[omProp,thick,dashed] table[x=t,y=maize]{figdata/markets-3/09-agriculturals-and-softs/mm.csv};\addlegendentry{maize price (right)}
\end{axis}
\end{tikzpicture}

\omcaption{CBOT corn: \omterm{def:m3:agriculturals-and-softs:cot}{managed money}'s net position as a share of open interest, monthly means of
weekly reports, January 2016 to July 2026, against the monthly maize price. Funds are long when prices
rise and short when they fall: the two series have a correlation of 0.68. Data: CFTC disaggregated
Commitments of Traders (futures only); IMF maize price via FRED.}\label{fig:m3:agriculturals-and-softs:cot}
\end{omfigure}

In the weekly corn data from January 2016 to September 2026, \omterm{def:m3:agriculturals-and-softs:cot}{managed money} was net long a record 24.4\%
of open interest on 22 March 2022 and net short 22.8\% on 9 July 2024; it was net short in 300 of 559
weeks, and on 15 September 2026 it was net long 22.5\%, 2.26 standard deviations above its average of the
previous three years. Traders read the report as a measure of crowding: a market in which funds are
already very long has fewer buyers left if the news disappoints.

\section{Softs and the crush}

\begin{definition}[Softs]\label{def:m3:agriculturals-and-softs:softs}
\emph{Softs}\index{softs} are agricultural commodities that are grown rather than mined and are not
grains or livestock: coffee, cocoa, sugar, cotton and orange juice.
\end{definition}

\omterm{def:m3:agriculturals-and-softs:softs}{Softs} are grown in few countries, often by smallholders, and their supply shocks are concentrated. The
IMF's monthly cocoa price averaged \$2{,}369 a tonne in 2022; it then rose more than fourfold, to a
monthly average of \$10{,}710 in January 2025, and was \$5{,}619 in July 2026
(\cref{fig:m3:agriculturals-and-softs:cocoa}). The International Cocoa Organization reported that both the 2022/23 and 2023/24
seasons ended in deficit, the second after weather, diseases and pests cut the crops of the main West African producers.

\begin{omfigure}
\begin{tikzpicture}
\begin{axis}[width=11cm,height=4.4cm,
  ylabel={thousand \$ per tonne},
  tick label style={font=\small},label style={font=\small},
  xmin=2000,xmax=2026.75,ymin=0,ymax=12,ytick={5,10},grid=major,grid style={gray!25},
  xtick={2000,2005,2010,2015,2020,2025},xticklabel style={/pgf/number format/1000 sep={}},
  table/col sep=comma]
\addplot[omProp,very thick] table[x=t,y=cocoa]{figdata/markets-3/09-agriculturals-and-softs/cocoa.csv};
\end{axis}
\end{tikzpicture}

\omcaption{Cocoa, monthly average price, January 2000 to July 2026. Data: IMF Primary Commodity Prices
(series PCOCOUSDM via FRED).}\label{fig:m3:agriculturals-and-softs:cocoa}
\end{omfigure}

Processing links crops to their products. A soybean crusher turns beans into meal (animal feed) and oil
(food and fuel).

\begin{definition}[Crush spread]\label{def:m3:agriculturals-and-softs:crush}
The \emph{crush spread}\index{crush spread} is the value of the meal and oil from one bushel of soybeans
minus the price of the bushel. At the exchange's conventional yields of 44 pounds of meal and 11 pounds of
oil per 60-pound bushel it is, in dollars per bushel,
$0.022 \times P_{\mathrm{meal}} + 0.11 \times P_{\mathrm{oil}} - P_{\mathrm{beans}}$, with meal in dollars
per short ton and oil in cents per pound.
\end{definition}

The crush is to a soybean processor what the crack is to a refiner
(\cref{ch:m3:refined-products-and-cracks}): its margin, hedged by buying beans and selling meal and oil,
in the ratio of ten soybean contracts to eleven of meal and nine of oil.

\section{Tutorial: positioning and limits}\label{tut:m3:agriculturals-and-softs:cot}

\begin{tutorial}
\textbf{Goal.} Turn the COT files into a positioning measure, and compute the limits and the lock of a
report day. \textbf{End state:}
\cref{fig:m3:agriculturals-and-softs:cot,fig:m3:agriculturals-and-softs:limits} and the margin of the
weekend problem.
\begin{tutsteps}
\item \textbf{Positioning.} Net positions, their share of open interest and a trailing z-score.
\omcode{code/firm/cot/firm_cot.py}{12}{32}{Net positions and their trailing z-score.}\label{lst:m3:agriculturals-and-softs:cot}
\item \textbf{Limits.} The variable-limit rule and the path of a locked market.
\omcode{code/firm/cot/firm_cot.py}{35}{63}{The limit rule and a limit-locked price path.}\label{lst:m3:agriculturals-and-softs:limits}
\item \textbf{Run} \texttt{m3\_ags.mm\_stats()}, \texttt{report\_day()} and \texttt{fig\_ags.py}.
\end{tutsteps}
\textbf{What to change next.} Compute the same measure for soybeans and wheat and see whether funds'
positions in the three move together; test whether extreme positioning predicts the next month's return,
correcting for the overlap of the weekly samples.
\end{tutorial}

\section{Build: the positioning parser and limit state}\label{bld:m3:agriculturals-and-softs:cot}

\begin{build}
\bfield{Purpose} The miniature firm trades grains around reports: it needs weekly positioning by
category, the current daily limits of each contract, and the state of a \omterm{def:m3:agriculturals-and-softs:limit}{limit-locked market} for its
risk system and its order router, which must not send orders that cannot trade.
\bfield{Interface} \texttt{net(row, category)}, \texttt{net\_share}, \texttt{zscore(values, window)};
\texttt{variable\_limit(average\_price)}; \texttt{limit\_path(start, fair, limit, expanded)};
\texttt{crush\_margin(beans, meal, oil\_cents)}.
\bfield{Rules} Positions are integers in contracts; z-scores use only past values; the \omterm{def:m3:agriculturals-and-softs:limit}{expanded limit}
applies only the day after a limit close; limits in dollars per bushel to the cent.
\bfield{Acceptance tests} \texttt{code/firm/cot/tests/}: net positions; the z-score window; the limit
rule at three price levels including the floor; a two-day lock; the board crush.
\bfield{Stretch} Feed limit states to \texttt{firm.calendar} (One Quant Book 1, chapter 22); read the
COT files directly from the Commission's archive format.
\end{build}

\begin{omsources}
\item US Commodity Futures Trading Commission, \emph{Commitments of Traders}: about the reports, release
schedule, disaggregated historical files 2016--2026.
\item USDA, \emph{WASDE}; National Agricultural Statistics Service, \emph{Acreage}, 30 June 2026; ERS
feed grains documentation (marketing years).
\item CME Group, CBOT rulebook chapter 10 (corn) and grain price-limit reset notices; soybean crush
reference guide.
\item International Monetary Fund, Primary Commodity Prices (cocoa, maize) via FRED.
\item International Cocoa Organization, \emph{Quarterly Bulletin of Cocoa Statistics}, November 2024 (summary).
\item CME Group, price limits page for agricultural futures, trade date 8 September 2026 (Internet Archive copy).
\end{omsources}

\section{Exercises}

\begin{exercise}[$\star$]\label{exo:m3:agriculturals-and-softs:1}
Ending stocks are 1.2 billion bushels and total use 14.8 billion. Give the \omterm{def:m3:agriculturals-and-softs:wasde}{stocks-to-use ratio}.
\end{exercise}

\begin{exercise}[$\star$]\label{exo:m3:agriculturals-and-softs:2}
Give the initial and \omterm{def:m3:agriculturals-and-softs:limit}{expanded limits} for corn averaging \$4.00 and \$6.00 a bushel.
\end{exercise}

\begin{exercise}[$\star$]\label{exo:m3:agriculturals-and-softs:3}
Why is July corn an old-crop contract and December a new-crop one?
\end{exercise}

\begin{exercise}[$\star\star$]\label{exo:m3:agriculturals-and-softs:4}
Meal is \$300 a short ton, oil 50 cents a pound, beans \$10.00 a bushel. Give the board crush.
\end{exercise}

\begin{exercise}[$\star\star$]\label{exo:m3:agriculturals-and-softs:5}
A report moves corn's fair value down \$0.60 with limits of 30 and 45 cents. How many days is a long
locked, and what does it lose per contract before it can sell?
\end{exercise}

\begin{exercise}[$\star\star$]\label{exo:m3:agriculturals-and-softs:6}
Funds hold a record net long just before a report that disappoints. Why might the fall be larger than
the report alone justifies?
\end{exercise}

\begin{exercise}[$\star\star\star$]\label{exo:m3:agriculturals-and-softs:7}
\emph{Coding.} With \texttt{mm\_series}, give \omterm{def:m3:agriculturals-and-softs:cot}{managed money}'s largest net long and net short in corn
since 2016, their dates, and the number of weeks net short.
\end{exercise}

\begin{exercise}[$\star\star\star$]\label{exo:m3:agriculturals-and-softs:8}
\emph{Find the flaw.} ``\omterm{def:m3:agriculturals-and-softs:cot}{Managed money}'s net position and the price move together, so the funds cause
the price moves; we will trade ahead of them.''
\end{exercise}

\section{Problem: Report Day}

\begin{problem}[{Weekend problem --- a long locked in by a report}]\label{pb:m3:agriculturals-and-softs:1}
A trader holds 50 corn futures long at \$4.50 a bushel. A report cuts the fair price to \$3.70. Limits
follow the rule of the dated box for an average price of \$4.50.

\medskip\noindent\textbf{Part I --- The limits.}
\begin{enumerate}
\item Give the initial and \omterm{def:m3:agriculturals-and-softs:limit}{expanded limits}.
\item What is the contract's value per bushel move?
\item Where does corn settle on the first day?
\item And on the second?
\item On which day can the trader sell, and at what price?
\end{enumerate}

\medskip\noindent\textbf{Part II --- The money.}
\begin{enumerate}[resume]
\item What variation margin does the trader pay on day 1?
\item And on day 2?
\item What margin must it hold to survive both locked days?
\item What is its total loss when it sells on day 3?
\item What would a stop-loss order at \$4.30 have done?
\end{enumerate}

\medskip\noindent\textbf{Part III --- Hedges.}
\begin{enumerate}[resume]
\item Could it have hedged with options during the lock? At what price?
\item Could it have sold a related contract (another month, another exchange)?
\item How does a producer hedging its crop experience the same days?
\item Why do exchanges keep limits despite this?
\item How would you size a position before a report, given the limits?
\end{enumerate}

\medskip\noindent\textbf{Part IV --- Judgement.}
\begin{enumerate}[resume]
\item What does the COT report tell you about the risk of such a report day?
\item Why is a limit not a cap on the loss?
\item How does a lock change the value of liquidity?
\item State the \emph{named result}: the margin needed to survive the lock, per contract and for the position.
\item In one sentence: what is a daily limit?
\end{enumerate}
\end{problem}

\section{Interview questions}

\begin{interviewq}[$\star$ \iqroles{trader}]\label{iq:m3:agriculturals-and-softs:1}
What is the \omterm{def:m3:agriculturals-and-softs:wasde}{stocks-to-use ratio} and why does the price depend on it non-linearly?
\end{interviewq}

\begin{interviewq}[$\star$ \iqroles{trader, researcher}]\label{iq:m3:agriculturals-and-softs:2}
What is a \omterm{def:m3:agriculturals-and-softs:limit}{limit-locked market}, and how do you manage the risk of one?
\end{interviewq}

\begin{interviewq}[$\star\star$ \iqroles{researcher}]\label{iq:m3:agriculturals-and-softs:3}
How would you build a signal from the COT report without look-ahead bias?
\end{interviewq}

\begin{interviewq}[$\star\star$ \iqroles{trader}]\label{iq:m3:agriculturals-and-softs:4}
Explain the soybean crush and how a processor hedges it.
\end{interviewq}

\begin{interviewq}[$\star\star$ \iqroles{risk}]\label{iq:m3:agriculturals-and-softs:5}
How do you compute a VaR-type measure for a position in a market with daily limits?
\end{interviewq}

\begin{interviewq}[$\star\star\star$ \iqroles{developer, trader}]\label{iq:m3:agriculturals-and-softs:6}
Design the report-day playbook system: what data, what checks, what automated actions at noon?
\end{interviewq}
